% !TEX program = xelatex
\documentclass[11pt,a4paper]{article}

\usepackage[UTF8]{ctex}
\usepackage{amsmath,amssymb,amsthm}
\usepackage{mathtools}
\usepackage{geometry}
\usepackage{hyperref}
\usepackage{xcolor}
\usepackage{enumitem}
\usepackage{booktabs}
\usepackage{longtable}
\usepackage{mdframed}
\usepackage{graphicx}
\usepackage{array}
\usepackage{bm}
\usepackage{caption}
\usepackage{tabularx}
\usepackage{float}

\geometry{margin=2.2cm, top=2cm, bottom=2cm}
\hypersetup{colorlinks=true, linkcolor=blue!60!black, urlcolor=blue!60!black}
\setlength{\parindent}{2em}
\setlist{nosep}

\definecolor{annotblue}{RGB}{46,168,229}
\definecolor{annotpink}{RGB}{220,50,180}

\newenvironment{important}{%
  \begin{mdframed}[linecolor=annotblue!70, backgroundcolor=annotblue!5, roundcorner=3pt]
  \textcolor{annotblue}{$\blacktriangleright$ \textbf{要点}}
  \medskip\par
}{\end{mdframed}}

\newenvironment{lesson}{%
  \begin{mdframed}[linecolor=annotpink!70, backgroundcolor=annotpink!5, roundcorner=3pt]
  \textcolor{annotpink}{$\blacktriangleright$ \textbf{教训}}
  \medskip\par
}{\end{mdframed}}

\title{\textbf{LLM Agent 驱动的论文数值复现}\\
  \large 以电磁多极展开（mie-f）四轮复现为例}

\author{张元正}
\date{2026-08-14}

\begin{document}

\maketitle
\thispagestyle{empty}

\begin{center}
  \fcolorbox{annotblue!50}{annotblue!5}{\parbox{0.92\textwidth}{\centering\small
  \textcolor{annotblue}{$\blacktriangleright$} \textbf{报告定位}：本文面向 LLM / agent 方向读者，总结我如何用多 agent 编排\\
  完成四篇电磁多极展开经典论文的数值复现。侧重 \textbf{agent 架构、工作流、记忆治理、验证方法论与经验沉淀}，\\
  不展开电磁学物理本身（物理推导与验证数字见另两份给光学组的报告）。}}
\end{center}

\tableofcontents
\newpage

\section{项目概览}

\subsection{复现对象}

本项目复现四篇电磁多极展开（multipole expansion）方向的经典论文，形成一条从「多极映射」到「精确多极矩」的理论链：

\begin{table}[H]
\centering\small
\begin{tabular}{@{}lll@{}}
\toprule
论文 & 内容 & 复现目标 \\
\midrule
Grahn 2012 (NJP 14 093033) & 电磁多极理论映射 & 电流矩 $\to$ 电/磁多极系数的映射关系 \\
FC2015 (Opt. Express 23 33044) & 精确偶极矩 & 精确电偶极矩定义与退化 \\
FC2017 (Sci. Rep. 7 7527) & toroidal 多极非独立 & toroidal 不是独立自由度的严格证明 \\
Alaee 2018 (Opt. Commun. 407 17) & 超越长波近似 & Table 1（长波）vs Table 2（精确）多极矩 \\
\bottomrule
\end{tabular}
\end{table}

\subsection{四轮复现与结果}

按「一轮复现一张图 / 一个结论」组织为四轮，每轮独立验收：

\begin{table}[H]
\centering\small
\begin{tabular}{@{}llll@{}}
\toprule
轮 & 对象 & 方法 & 结果 \\
\midrule
1 & Fig.1 介电球散射截面 & Mie 解析 + Table 2 积分 & 双实现交叉 $<0.21\%$ \\
2 & Fig.2 金球多极分解 & 同上（加密网格） & 双实现交叉 $<0.03\%$ \\
3 & Grahn 映射 & 双路径 + 独立远场 & 相对差 $5.7\times10^{-7}\sim2.9\times10^{-4}$ \\
4 & Fig.3 双金盘 Fano & COMSOL FEM + BEM & 两求解器交叉，网格收敛评估 \\
\bottomrule
\end{tabular}
\end{table}

\begin{important}
一个关键数字：双金盘（非球形、无解析解）用边界元法（BEM）$128$ GB 就跑通了同一几何，而有限元（FEM）频域直接求解器需要 $900$ GB——省约 $7$ 倍。这说明**选对数值方法比堆算力更重要**。
\end{important}

\subsection{本报告关注的维度}

给光学组的报告讲「物理为什么对」（三条物理判据），本报告讲「agent 怎么把这件事做对」——即复现过程中 \textbf{agent 编排、记忆治理、验证方法论、读图处理、经验沉淀}这一整套工程方法。二者互补，物理报告回答「算对了吗」，本报告回答「这套 agent 工作流能不能复用」。

\section{Agent 架构与分工}

复现不是单模型从头做到尾，而是**三层 agent 分工**——每层有不同的职责、模型档位与上下文边界。

\subsection{三层分工}

\begin{table}[H]
\centering\small
\begin{tabular}{@{}p{2.6cm}p{4.6cm}p{2.2cm}p{2.2cm}@{}}
\toprule
层 & 职责 & 模型档位 & 上下文边界 \\
\midrule
\textbf{主 agent} & 编排：读路线图、spawn 子 agent、汇总报告、在 gate 停顿问用户 & 主循环 & 保留判断层，不亲自做隔离活 \\
\textbf{子 agent} & 执行单步（formalization / 实现 / 验证…） & 常规档 & 只读/写限定目录，报告回传结论 \\
\textbf{codex 后台委托} & 多模态读图、并发批量、长任务 & 难活用强模型 & 后台 bash 并发，产物落盘不占主 context \\
\bottomrule
\end{tabular}
\end{table}

\subsection{核心纪律：判断密度路由}

一条贯穿始终的规则——\textbf{读文件内容不进主 context 的，都派子 agent}。主 agent 亲读大段文件/多文件会把原文灌进主 context，压缩信噪比。所以：

\begin{itemize}
  \item 单事实查询（已知文件/符号）→ 直接 Grep/Read，不派 agent；
  \item 跨多文件调查、架构梳理 → 派只读子 agent，只收结论不收文件转储；
  \item 裁决必需的读（gate 停机点）→ 主 agent 亲读，不转述。
\end{itemize}

\subsection{子 agent 并发（fan-out）}

两张独立图 / 两个独立子任务可并发 spawn 多个子 agent，主 agent 是唯一汇聚点。\textbf{并发前提是真独立}——无数据/文件/逻辑依赖；有依赖就串行。

\subsection{codex 后台委托（并发 + 长任务）}

多模态读图、大批量、深推理的任务交给 codex 后台执行，特点是：

\begin{itemize}
  \item \textbf{并发 = 多条后台 bash}，每条 \texttt{run\_in\_background}，互不阻塞，各自落盘；
  \item \textbf{长任务不设硬超时}——委托可能跑很久，硬超时会误杀接近完成的任务（踩过：codex 渲染完大部分产物却在收尾落盘时被超时 kill）；
  \item \textbf{判活看落盘产物}是否在增长（事件流尾部时间戳、result 文件是否出现），而非猜；
  \item \textbf{产物落盘传路径}，不靠 stdout（会截断），主 agent 亲读落盘产物做验收。
\end{itemize}

\begin{important}
这套三层分工的关键不是「用多个模型」，而是**把「判断」和「执行」分离**：判断（gate 裁决、物理归因、方法选型）留在主循环，执行（写码、跑脚本、读图、批量）下沉到子 agent / 后台。这样主 context 始终只装「结论 + 关键定位」，而不是执行细节。
\end{important}

\section{复现工作流}

\subsection{多轮 × 7 步循环}

每轮复现一张图/一个结论，轮内 7 步：

\begin{table}[H]
\centering\small
\begin{tabular}{@{}lll@{}}
\toprule
步 & 名 & 产出 \\
\midrule
01 & PDF 预处理 & 从论文提取参数/公式/图数据 \\
02 & formalization & 参数确认 + 机器可判 spec（YAML） \\
03 & theory notes & 公式来源、推导、与教材对标 \\
04 & implementation & 代码 + 测试（TDD，物理约束先硬编码） \\
05 & run \& monitor & 运行 + 收集数据 \\
06 & physical verification & 3 层验证（硬约束→极限→论文图量化） \\
07 & analysis \& report & 归因 + 双报告 + 记忆 \\
\bottomrule
\end{tabular}
\end{table}

轮间共享公共模块（特殊函数、基准解、多极矩），每轮开头先读上一轮 worklog 决定复用 vs 新建。

\subsection{4 个人工 gate（agent 自由跑，gate 必须停）}

agent 可以自动推进，但四个节点必须停下来问用户，不代替人工裁决：

\begin{table}[H]
\centering\small
\begin{tabular}{@{}lll@{}}
\toprule
gate & 触发点 & 用户核对什么 \\
\midrule
① 参数 gate & formalization 末 & 参数、单位、范围 \\
② spec gate & formalization 末 & 物理形式化与论文一致 \\
③ 公式 gate & 实现/运行末 & 核心公式对权威教材核 \\
④ 误差 gate & 报告末 & 量化误差数字，不接受「看起来一致」 \\
\bottomrule
\end{tabular}
\end{table}

\subsection{复述纪律（防转述漂移）}

主 agent 向用户汇报任何「gate 裁决 / verifier 输出 / 已归档结论」的量化数值时，\textbf{必须现场重开原始文件核对}，不得凭对话历史记忆转述。数字、方向词、范围必须与原文逐字一致。这一条专门防「模型在长对话里凭印象转述已裁决数字，结果越传越偏」。

\subsection{人工复核硬化}

凡从论文图/位图提取的参数（坐标轴类型、刻度值、峰位、曲线数据），模型读数一律视为「未核实线索」：

\begin{enumerate}
  \item 下一步开工前必须人工复核——formalization 定稿前用户亲自核对；
  \item workflow 结束前整体复核一次——gate④ 前把所有「从图提取的参数」列清单，用户逐项复核再定稿 REPORT。
\end{enumerate}

\subsection{关键节点必须停}

以下节点默认停（除非用户说全自动）：进 REPORT 前、改 skill/蓝图前、物理验证失败要换方案时、遇到缺失信息时、\textbf{偏离既定 workflow 步骤时}。最后一条尤其重要——不得自主偏离后在报告里事后声明代价，偏离前必须问用户。

\section{记忆与知识管理}

复现跨多轮、多会话，agent 需要四层记忆载体，各司其职、不混用：

\begin{table}[H]
\centering\small
\begin{tabular}{@{}p{3.2cm}p{3.4cm}p{3.4cm}p{2.6cm}@{}}
\toprule
载体 & 职责 & 时效 & 写入时机 \\
\midrule
\textbf{memento} & 跨会话长期：事实/决策/坑，语义检索 & 永久，按 importance 分层 & 有实质产出/踩坑/决策时 \\
\textbf{MEMORY.md} & 恢复上下文：「我现在在哪、要做什么」 & 短时态，覆盖式更新 & 每工作单元结束 / compact 前 \\
\textbf{doc-sync} & 面向用户的设计现状 + 待决策 & 长期 & 新 agent 完成 / 新决策时 \\
\textbf{skill} & 可复用指令：经验沉淀为 SKILL.md & 永久，跨场景 & 经验反复复现、用户拍板后 \\
\bottomrule
\end{tabular}
\end{table}

\subsection{memento：语义检索的长期记忆}

用「关键词 + 向量」融合排序，支持跨会话召回「三年前踩过的坑」。关键做法：

\begin{itemize}
  \item \textbf{存完整句子而非碎片}：内容越完整向量语义越准；
  \item \textbf{importance 分层}：决策 0.7+、日常 0.3--0.5，避免低价值记忆淹没搜索；
  \item \textbf{存前查重}：用向量余弦相似度检查是否已有相似记忆；
  \item \textbf{结构化三工具}：事实/决策走 \texttt{memory\_store}，架构决策走 \texttt{decisions\_log}，踩坑走 \texttt{pitfalls\_log}。
\end{itemize}

\subsection{MEMORY.md：恢复上下文的索引头}

一个纯索引头文件（$\le200$ 行），会话启动注入，回答「我现在在哪、当前在做什么」。\textbf{只放索引，不放正文}——发现、行为规范、任务细节全部下沉到 \texttt{memory/} 目录的 topic 文件，头部只留指针。每次 compact 前覆盖式更新，不累积瞬时状态。

\subsection{doc-sync：面向用户的设计同步}

一套「主文档 + 副文档 + 待决策清单」系统，作为唯一 human 入口。用户只看 \texttt{00-待决策清单} 就能拍板，其余读完即可。主文档多用 ASCII 图表示流程 + 状态表 + 跳转链接，副文档放每个设计 flow 的详细设计。

\subsection{skill：从经验到可复用指令}

复现结束后，把反复用到的经验提炼为 skill（本报告 §7 详述 10 个 skill）。核心原则是\textbf{「泛化而非复制」}——把项目专属的细节剥离，只保留可跨任务复用的编排/方法/铁律。

\begin{important}
四层载体的边界判据：\textbf{需要语义召回 + 结构化} → memento；\textbf{恢复「我在哪」} → MEMORY.md；\textbf{给用户看现状与待决策} → doc-sync；\textbf{跨场景稳定复用} → skill。混用会各自失效——例如把「当前在做什么」塞进 memento，下次召回时反而找不到「长期结论」。
\end{important}

\section{验证方法论}

复现的落脚点不是「公式算对了」，而是「凭什么相信自己算对了」。这需要一套机器可判的验证方法。

\subsection{formalization：把论文声明规格化}

第一步把论文的图/表/结论翻译成机器可判的 YAML spec——参数、单位、范围、验证判据全部显式写死。这是后续所有验证的契约：agent 对着 spec 跑，而不是对着「看起来像不像」跑。

\subsection{物理验证三元组}

三条\textbf{相互独立}的判据，任何一条不过都不放过：

\begin{enumerate}
  \item \textbf{多近似方法对比}（排除实现 bug）——同一物理对象用两条独立数学路径算，交叉比对。球体：Table 2 体积积分 vs Mie 级数求和；非球形无解析解：两种数值求解器互验（FEM vs BEM）。
  \item \textbf{极限退化}（排除公式抄错）——每个公式要求它在已知极限退化回已知结果。精确 Table 2 在 $kr\to0$ 退化回长波 Table 1；Mie 退化回 Rayleigh $x^4$；$Q_{\rm ext}\to2$。多极矩系数由退化固定，不是照抄论文。
  \item \textbf{能量守恒 / 光学定理}（排除物理不自洽）——$C_{\rm ext}=C_{\rm sca}+C_{\rm abs}$ 到 $10^{-10}$；光学定理 $S(0)$ 双路差 $10^{-14}$；功率闭合 balance $=1.0$。
\end{enumerate}

三条叠加排除「实现 bug」「公式抄错」「物理不自洽」三类错误。

\subsection{预注册：先冻结 spec 再跑}

验证阈值（如 RMSE 0.020、p95 0.050）\textbf{必须在看到结果前冻结}，防「看到结果后调阈值抹平差异」这种事后发明。这是典型的科学诚信失败模式，agent 尤其容易踩——所以要显式预注册。

\subsection{result\_class 诚实标注}

结果用 7 级枚举标注，\textbf{禁把「跑通了管道」当「复现成功」}：

\begin{table}[H]
\centering\small
\begin{tabular}{@{}ll@{}}
\toprule
等级 & 含义 \\
\midrule
\texttt{physical\_reproduction\_success} & 物理复现成功 \\
\texttt{partial\_physical\_match} & 部分物理吻合（缺网格收敛/远场闭合等） \\
\texttt{surrogate\_fallback} & 替代方案（如读图数据） \\
\texttt{diagnostic\_only} & 仅诊断性结果 \\
\texttt{pipeline\_completed} & 管道跑通，但物理未验证 \\
\texttt{not\_run} / \texttt{failed} & 未运行 / 失败 \\
\bottomrule
\end{tabular}
\end{table}

读图数据永远是 \texttt{surrogate\_fallback}，不是物理复现。「跑通了管道」（\texttt{pipeline\_completed}）≠ 物理结论成立。

\section{读图处理与人工复核}

\subsection{为什么读图是最大隐性错误源}

复现论文离不开从图里读参数（坐标轴类型、刻度值、峰位、曲线数据）。但模型读图极不可靠——\textbf{模型看同一张图两次可能给不同答案}。

\begin{lesson}
最典型的一次：Fig.1 的 y 轴被模型判为 log 坐标，实际是 linear（刻度 1,3,5,7）。这个误判让整轮结论报废，最后被 vision-mcp 4 个通道交叉（2 个裁剪图 + 整图独立读）推翻。教训：\textbf{任何单通道读数（包括主 agent 亲看原图）都可能错}。
\end{lesson}

\subsection{刻度/坐标轴铁律}

从论文图提取任何参数时，必须遵守：

\begin{enumerate}
  \item 用权威最高分辨率原图，不用低清 OCR 渲染、不用缩略图；
  \item \textbf{vision-mcp 多通道收敛}：$\ge2$ 裁剪图 + 整图独立读，全一致才可信；
  \item \textbf{轴标题 vs 刻度分离}：竖排英文轴标题（如 \texttt{(C\_sca/(}$\lambda^2$\texttt{/2}$\pi$\texttt{))}）是文字不是刻度，先剥离标题区域，否则竖排笔画混入被误判为刻度线；
  \item 读到的值标 confidence + 提取方法 + 「未核实线索，待人工复核」。
\end{enumerate}

\subsection{log / linear 判定}

不凭直觉判 log/linear。判定信号：刻度间距是否等距（linear 等距，log 每十倍等距）、轴标签是否含 10 的幂、网格线分布。多通道读同一根轴全一致才算定论，不一致就标「未核实，待复核」，\textbf{不硬选一个继续算}。

\subsection{读图数据 vs 解析解的取舍}

纯理论论文（曲线 = 标准解析解，如 Alaee 2018 的曲线就是标准 Mie）\textbf{没有作者原始数据可对齐}，直接自己算解析解，读图数据放弃，人工目视相似即可。读图数据只用于峰位/形态定性比对，不用于精确数值标定——它永远是 \texttt{surrogate\_fallback}。

\section{挑战、教训与 skill 沉淀}

\subsection{踩过的主要坑}

\begin{table}[H]
\centering\small
\begin{tabular}{@{}p{4.2cm}p{9.0cm}@{}}
\toprule
坑 & 教训 \\
\midrule
codex 反复被杀 & 服务端过载时后台任务被中断；resume 对多 agent 子会话不可用，只能 fork 重发；单任务连断多次 = 停发等窗口 \\
长任务硬超时 & 委托接近完成时被超时 kill；改为不设超时，靠结束唤醒 + 落盘产物判活 \\
读图误判 & log/linear 误判毁掉整轮；必须 vision-mcp 多通道收敛 + 人工复核 \\
BEM 调试 15 轮 & 接口不熟导致反复试错；核心坑：色散金属域 BEM 不支持、表面网格 selection 维度、远场需 FarFieldCalculation \\
内存估算 & FEM 直接求解器 $O(\mathrm{DOF}^{1.5})$，0.3 网格需 900G；先估 DOF 再申请资源 \\
上下文膨胀 & 主 agent 亲读大文件撑爆 context；判断密度路由，读文件不进主 context 的派子 agent \\
\bottomrule
\end{tabular}
\end{table}

\subsection{从经验到 skill：10 个可复用 skill}

复现结束后，把反复用到的经验提炼为 skill。共沉淀 \textbf{10 个}——其中 6 个从本项目直接提炼，4 个是过程中用到的通用能力：

\begin{table}[H]
\centering\small
\begin{tabular}{@{}p{4.6cm}p{7.6cm}@{}}
\toprule
skill & 内容 \\
\midrule
\textbf{paper-reproduction} & 复现编排：多轮×7步 + 4 gate + 复述纪律 + 人工复核 \\
\textbf{figure-reading} & 读图提取：刻度铁律 + 多通道收敛 + 数字化 \\
\textbf{comsol-scattering} & COMSOL FEM/BEM 建模 + 多极矩后处理 + 网格收敛 \\
\textbf{magnus-submit} & 集群 blueprint 提交 + 资源声明 + 红线 \\
\textbf{numeric-verification} & 物理验证三元组 + Richardson + result\_class \\
\textbf{third-party-cross-validation} & 第三方交叉验证 + 诚实区分同/异对象 \\
\midrule
adversarial-review & 对抗性审查（独立视角审代码/结论） \\
transition-wrapup & 会话收尾（整理 memento/MEMORY/rules） \\
doc-sync & 设计文档同步（面向用户的现状+待决策） \\
grill-me & 追问（对 agent 结论追问拍板） \\
\bottomrule
\end{tabular}
\end{table}

\subsection{skill 提炼原则：泛化而非复制}

把一个项目专属的经验变成 skill，关键动作是\textbf{剥离专属细节、保留可复用骨架}：

\begin{itemize}
  \item 项目内的 \texttt{agent-workflow}（写死「4 轮、Alaee Fig.1/2/3」）$\to$ 泛化为 \texttt{paper-reproduction}（「多轮、任意论文」）；
  \item 散在多个 skill 里的「刻度铁律」$\to$ 收口成独立的 \texttt{figure-reading}（单一权威来源）；
  \item 物理细节（Mie 公式、具体参数）不进 skill，留在各轮自己的 \texttt{formalization/}。
\end{itemize}

\begin{important}
skill 的边界判断：\textbf{「脱离本项目，这件事能否独立成立为一个可复用任务？」}能 → 独立 skill；不能（只是某项目的局部步骤）→ 留在项目内。避免把「我们做过 X」都堆成 skill——碎片化的 skill 比没有更糟。
\end{important}

\section{结论与可复用性}

\subsection{这套方法沉淀了什么}

四轮复现沉淀了一套\textbf{agent 驱动的科研计算方法论}，核心是五个可复用组件：

\begin{enumerate}
  \item \textbf{三层 agent 分工}——判断与执行分离，主 context 只装结论；
  \item \textbf{gate + 复述纪律 + 人工复核}——把「可信」交给机器可判的契约 + 人工停点，不交给模型自觉；
  \item \textbf{四层记忆载体}——memento / MEMORY.md / doc-sync / skill 各司其职；
  \item \textbf{物理验证三元组 + result\_class}——回答「凭什么相信算对了」并诚实标注完成度；
  \item \textbf{skill 沉淀}——10 个可复用 skill，经验不随会话消失。
\end{enumerate}

\subsection{关键数字一览}

\begin{table}[H]
\centering\small
\begin{tabular}{@{}ll@{}}
\toprule
指标 & 值 \\
\midrule
复现论文 & 4 篇，4 轮 \\
球体双实现交叉 & $<0.21\%$（介电）/ $<0.03\%$（金） \\
Grahn 映射相对差 & $5.7\times10^{-7}\sim2.9\times10^{-4}$ \\
能量守恒 & 到 $10^{-10}$ \\
光学定理双路差 & $10^{-14}$ \\
BEM vs FEM 内存 & $128$G vs $900$G（省 7 倍） \\
沉淀 skill & 10 个 \\
\bottomrule
\end{tabular}
\end{table}

\subsection{迁移价值}

这套方法不限于电磁散射，对任何「有解析基准可校验、有数值可交叉」的科研计算复现都成立：

\begin{itemize}
  \item \textbf{有解析解的}（球体/Mie 类）→ 双实现交叉 + 极限退化，解析解就是验金石；
  \item \textbf{无解析解的}（非球形/数值类）→ 先球体校验求解器，再两数值方法互验 + 网格收敛；
  \item 无论哪类，\textbf{formalization 规格化 + 三元组 + result\_class 诚实}都是通用底座。
\end{itemize}

\begin{important}
一句话总结：\textbf{复现的可信度不是「看起来像」，而是「机器可判的契约 + 相互独立的物理判据 + 诚实标注的完成度」。} agent 的价值在于把这三件事制度化、可复用，而不是替人「画一张像的图」。
\end{important}


\end{document}
